\section{One diagonal power serves three purposes}\label{sec:diagonal}

The diagonal choice $F=t^t$, the alias $R=F$, and the linear odd beta
spacing already occur in the incoming Linear CRT manuscript~\cite{linearcrt}.
We give the needed proofs here to make the new sign-bootstrap and pooling
composition reviewable. The distinct feature of this source is the fixed
radix class $q\equiv3\pmod6$.

The direct arithmetic construction encodes a finite multiset of valid
instances of a single-conclusion proof calculus. An exact balance between
conclusion labels and premise labels characterizes a proof of the target.
The existing compiler represents a finite sequence of multiplicative
accumulator states through beta remainders and certifies its local rows
through amplified binomial CRT conditions. We retain that architecture
and change the growth parameter and sign bootstrap.

Throughout this section, $t\ge2$ is an even integer. In the actual
polynomial source, $t=16+2h^2$, so these premises are automatic. Put
\[
 F=t^t,\qquad L_t=\lcm(1,2,\ldots,t).
\]
The variable exponent is proof notation: $F=t^t$ is enforced by an
existing simultaneous Pell power certificate in the ordinary source.
It is not an operation in the final quartic.

\subsection{Separating bounded multisets at one integer}

For a finite list $A=(a_1,\ldots,a_m)$ with $1\le a_i\le t$ and
$0\le m\le t$, define the auxiliary polynomial
\[
 H_A(X)=\prod_{i=1}^{m}(1+a_iX).
\]
The empty product is one. This polynomial is used in the proof only;
the length of the list is not the number of final unknowns.

\begin{lemma}[Diagonal evaluation separates the lists]\label{lem:labels}
If $A$ is such a list, $B$ is any finite list of labels in
$\{1,\ldots,t\}$, and $H_A(F)=H_B(F)$ with $F=t^t$, then $A$ and $B$
have the same multiplicities of every label. Only one initial length
bound is required.
\end{lemma}
\begin{proof}
First their lengths are equal. If $m>n$, then $m\ge n+1$ and
\[
 H_A(F)\ge(F+1)^m\ge(F+1)^{n+1}
       >(t(F+1))^n\ge(tF+1)^n\ge H_B(F).
\]
The strict inequality follows from $F+1>t^n$. The following weak
inequality also covers $n=0$. If $n>m$, the symmetric comparison uses
$m\le t$ and $F+1>t^m$; no prior upper bound on $n$ is needed.

Let the common length be $m$. For $0\le k\le m$, the coefficient of
$X^k$ in either product is a positive integer and is at most
\[
 \binom{m}{k}t^k\le\binom{t}{k}t^k
 \le t^{t-k}t^k=t^t=F.
\]
Here $\binom{t}{k}=\binom{t}{t-k}\le t^{t-k}$.
Thus every coefficient of $H_A-H_B$ has absolute value strictly less
than $F$. If a nonzero coefficient had least index $k$, division of
$(H_A-H_B)(F)=0$ by $F^k$ and reduction modulo $F$ would force that
coefficient to be divisible by $F$, a contradiction. Therefore
$H_A=H_B$ in $\Z[X]$.

Each $1+jX$ is a primitive linear polynomial, and two such factors with
positive $j$ are associates only when their labels agree. Unique
factorization, or comparison of the rational roots $-1/j$, now gives
the asserted multiplicities.
\end{proof}

Truth label zero has weight one and is deliberately omitted from these
lists. The retained proof-flow interface treats truth as freely available.
The lemma therefore provides exactly the projected label balance needed
by that interface. It does not claim that all integers $1+jF$ are pairwise
coprime; they need not be.

\subsection{The same power covers the CRT rows}

\begin{lemma}[Elementary least-common-multiple bound]\label{lem:lcm}
For even $t\ge2$, $L_t^2\le t^t=F$.
\end{lemma}
\begin{proof}
Write $L_t$ as the product of the largest power of each prime not
exceeding $t$. Each factor is at most $t$. There are at most $t/2$
such primes: besides 2, they belong to the $t/2-1$ odd integers from
3 through $t-1$. Thus $L_t\le t^{t/2}$, and squaring proves the claim.
\end{proof}

The inherited factorial-free CRT theorem uses
\[
 R\ge L_t^2,\qquad B=\binom{R+t}{t}.
\]
It proves two different prime-power facts. Every row $i\in\{1,\ldots,t\}$
has a full prime power $p^e\Vert B$ with $p^e\mid R+i$ and $p^e>t$;
conversely every full prime power of $B$ divides at least one such row.
The first fact supplies row-by-row soundness. The second supplies a
consistent reverse CRT construction at all prime factors of $B$.
The retained proof includes their valuation arguments and the amplified
range extraction; this report does not replace them by an assumption of
pairwise coprimality of the consecutive factors.

Lemma~\ref{lem:lcm} permits the literal specialization
\begin{equation}\label{eq:shift}
 R=F=t^t,\qquad B=\binom{F+t}{t}.
\end{equation}
The former square relation $R=F^2$ is unnecessary. This is a change to
the semantically certified source, followed by a polynomial alias and
exact elimination of its now redundant coordinate. It is not the
unjustified deletion of an equation from the old zero set.

\subsection{The required parity is already present}

The same choice also provides
\[
 v_2(F)=t\,v_2(t)\ge t>v_2(t!).
\]
Consequently an arbitrarily large odd multiple of $F$ can be chosen
divisible by $t!$: multiply by an odd multiple of the odd part of $t!$.
This is precisely the freedom used to choose a beta spacing large enough
for every pending proof-accumulator state. The polynomial need not
represent $t!$; it occurs only in the constructive existence proof of an
appropriate spacing.

Also $B=\binom{F+t}{t}$ is odd. For $1\le i\le t$,
$v_2(i)<v_2(F)$, so $v_2(F+i)=v_2(i)$. The powers of two in the
numerator product and in $t!$ therefore cancel exactly. This parity
allows the subsequent CRT representatives and beta codes to be chosen in
the required residue classes.

\section{Recovering signs before arithmetic interpretation}\label{sec:signs}

The direct final equation has unrestricted integer unknowns. Positivity
must therefore follow from equations, and the order in which it follows
matters. We replace the old positive-square radix by
\begin{equation}\label{eq:radix}
 q=6z+3,\qquad H=3+6\sum_{j=1}^{9}a_j,\qquad q=(BH)^x,
\end{equation}
where the $a_j$ are the nine actual CRT representatives. The last identity
is an existing polynomially certified power request.

Before the sign of any representative or of $B$ is known,
\[
 |q|\ge3,\qquad |q+1|\ge2.
\]
Thus the signed power theorem applies to both bases. The requests with
base $q+1$ force the nine representative indices to be nonnegative;
the request with base $q$ and index $y+1$ gives $y+1\ge0$.
It follows that $H\ge3$ and $y\ge-1$ without assuming a binomial value
or a positive radix.

The beta spacing can now be linear in its cofinal parameter:
\begin{equation}\label{eq:beta}
 y=F(2v+1).
\end{equation}
Because $F=t^t\ge2$ and $y\ge-1$, its odd multiplier cannot be negative.
Hence $y\ge F$ and $v\ge0$. This recovers the sign that a square formerly
imposed, while retaining a cofinal collection of legal spacings. The full
bootstrap, including the normalization ratios, the signed first binomial
digit, $B>0$, $x>0$, and finally $q>0$, is given in the following proof.

\subsection{The signed power and remainder interfaces}

The unchanged semantic anchor and full orientation frame supply
$W\ge258$, $|A|>W^W$, canonical signed Pell pairs, and bounds on the
actual power bases, indices and outputs. The two pooled blocks are
covered by the theorem just proved. For the three individual requests,
the usual congruence
\[
 X_k(A)+(n-A)Y_k(A)\equiv m\pmod{2An-n^2-1}
\]
gives the same conclusion when $|n|\ge2$: the bounded congruence
is the equality $m=n^k$, and a negative $k=-j$ would force
$mn^j=1$. This is impossible for an integer $m$. The bounds use
absolute values and do not require $n$ to be positive.

The common normalized remainder certificate is
\[
 M^2=1+4\sum_j R_j^2+a^2+b^2+c(c+1).
\]
As $c(c+1)\ge0$ for every integer $c$, one has $M\ne0$ and
$|R_j|<|M|/2$. If $L_jv_j=M$ and $v_jU_j=MQ_j+R_j$, then
$L_j,v_j\ne0$, $v_j\mid R_j$, and the integer $r_j=R_j/v_j$ satisfies
\[
 U_j=L_jQ_j+r_j,\qquad |r_j|<|L_j|/2.
\]
A congruent remainder in this open centered interval is unique, even
for negative $L_j$. The first binomial interpretation can therefore be
recovered before the radix has been proved positive.

\subsection{An acyclic bootstrap}

The horizon $t=16+2h_0^2$ is positive and even. Its individual power
request first gives $F=t^t\ge2^{64}$.
The fixed class $q=6z+3$ makes the two pooled bases nonzero nonunits.
Their power requests give $x\ge0$, $y+1\ge0$, and nonnegative values
of all nine actual representatives $a_j$. Hence $H\ge3$.
From $y=F(2v+1)$ and $y\ge-1$ follows $y\ge F$.

The raw index equation is
\[
 B\,u=C_0+F+1,
\]
where $C_0$ is the already nonnegative actual index representative.
Thus $B\ne0$ and $|BH|\ge3$. The individual self-radix request
gives $q=(BH)^x$. Since $q\ne1$, one has $x\ge1$.
The first and last beta endpoint equations are
\[
 x=F+1+(y+1)u_0,
 \qquad x=1+\bigl(1+(t+1)y\bigr)u_t.
\]
If $u_0\le-1$, the first gives $x\le F-y\le0$, impossible.
Therefore $x\ge F+1$. The last endpoint then forces $u_t\ge1$, so
\begin{equation}\label{eq:xlarge}
 x\ge2+(t+1)y\ge2+(t+1)F.
\end{equation}
Write $Q=|q|$ and $b=|B|$. Without assuming either sign, we now have
\[
 Q=(bH)^x\ge2^x,\qquad Q\ge9b^2,\qquad b\le Q/9.
\]
All low-digit denominator exponents are nonnegative and at most $x$.
Their exact ratios to the common modulus $M=q^x$ are integral.

\subsection{The first binomial digit forces a positive radix}

For $0\le j\le t$, put $C_j=\binom{F+t}{j}$.
Since $F\ge t$, these coefficients increase up to $t$, and
\[
 C_j\le C_t=\prod_{i=1}^t\left(1+\frac{F}{i}\right)
       \le(F+1)^t=:V_*.
\]
The bound $F+1\le2^{F-1}$ and \eqref{eq:xlarge} imply
\[
 \frac{V_*}{Q}\le2^{-F-t-2},\qquad \frac{V_*}{Q-1}<1/32.
\]
The authenticated numerator is $U=(q+1)^{F+t}$, and its digit row is
\[
 U=q^t(B+q\xi)+r,\qquad |r|<Q^t/2.
\]
For either sign of $q$, the binomial theorem gives the canonical
integer remainder
\[
 r_0=\sum_{j=0}^{t-1}C_jq^j,
 \qquad |r_0|<\frac{V_*Q^t}{Q-1}<Q^t/32.
\]
Centered remainder uniqueness yields $r=r_0$, and division by $q^t$
gives $B\equiv C_t\pmod q$. But $|B|\le Q/9$ and
$C_t<Q/32$, so $|B-C_t|<Q$. Thus
\[
 B=\binom{F+t}{t}>0,\qquad q=(BH)^x>0,\qquad N=B^x>0.
\]
Because $t\ge2$, this binomial coefficient exceeds $F+1$.
The earlier divisibility $B\mid C_0+F+1$ excludes $C_0=0$;
hence $C_0\ge1$, as required by the row-index extraction.
This argument handles the entire negative-radix possibility before
excluding it.

\subsection{The two prime-power covering statements}

For completeness, let $\Lambda=\lcm(1,\ldots,t)$, $R\ge\Lambda^2$,
and $B=\binom{R+t}{t}$. Every row $i\in\{1,\ldots,t\}$ has a full
prime power $p^e\Vert B$ dividing $R+i$ with $p^e>t$.
Indeed, if every $v_p(R+i)\le2\lfloor\log_p t\rfloor$, then
$R+i\mid\Lambda^2$, contradicting $R+i>\Lambda^2$.
Choose $p$ with $v=v_p(R+i)>2a$, $a=\lfloor\log_p t\rfloor$.
For $j\ne i$, the distance $|j-i|\le t$ implies
$v_p(R+j)=v_p(j-i)$. Therefore
\[
 e=v_p(B)=v-v_p\!\left(t\binom{t-1}{i-1}\right).
\]
In the Legendre sum for the subtracted valuation, each level up to $a$
contributes zero or one, and all later levels contribute zero.
Thus $v-a\le e\le v$, which proves $p^e>t$ and $p^e\mid R+i$.

Conversely, every full prime power of $B$ divides at least one row.
For a fixed prime let $v=\max_i v_p(R+i)$. Legendre's formula for
$v_p(B)$ is a sum of zero-or-one interval-carry terms, with no terms
above level $v$. Hence $v_p(B)\le v$. Both statements are needed:
one supplies a large row selector, and the other covers every prime
factor during reverse CRT construction. Consecutive $R+i$ need not
be pairwise coprime.

Here $R=F=t^t\ge\Lambda^2$ by Lemma~\ref{lem:lcm}, so the
retained factorial-free binomial CRT compiler applies.

\subsection{Remaining local bounds and the beta endpoint}

The inherited local-polynomial ledger uses exponent $2^{23}$.
For $R_{\rm maj}=2+t+x+y+F$, the established bounds give
$R_{\rm maj}<2x$ and $x\ge2^{64}+1$. The retained uniform estimate
$(2x)^{2^{23}}<2^x$ therefore applies. Its derived-label degree
bound $2^{22}(t+1)<x$ also holds. Since $H=3+6\sum a_j$ dominates
the old representative majorant, the other ten binomial digit
remainders retain their strict relative margin below $1/32$.
Their exact divisibilities give the same local polynomial constraints,
with the strengthened strict local ranges from the pooling source.

For each row the selected prime power $p^e>t$ identifies the correct
index. The amplified range lemma, the precision $x$, and the bound
$0\le L<2^x\le p^{ex}$ turn the local congruence into $L=0$.
Thus all supplied Horn-rule and accumulator rows are genuine.

There is a small boundary case to settle. If $y=F$, the first beta
endpoint gives initial proper remainder zero, while the last gives
final remainder one. The recovered local equations are
\[
 S_{i+1}W_{\mathrm{concl}(i)}
  =S_iW_{\mathrm{prem}_1(i)}W_{\mathrm{prem}_2(i)},
 \qquad W_j=1+jF>0.
\]
A zero initial state propagates to every state, contradicting the last.
Hence $y>F$. The actual first remainder is $F+1$, the last is one,
and all intervening states are positive. Their product identity and
Lemma~\ref{lem:labels} recover the exact multiplicities of proof labels.
The retained single-conclusion flow criterion then yields a $C_{222}$
refutation. It allows cycles and repeated premises: summing the balance
over labels outside the finite closure of truth gives the usual
contradiction if the target is outside that closure.

\subsection{Constructive completeness}

From a finite $C_{222}$ refutation, choose a finite proof tree and an
even horizon $t=16+2h_0^2$ strictly larger than its size, labels and
every local parameter. Pad with truth to $t$ rows. Set $F=t^t$ and
form the finite pending-multiset accumulator path with weights $1+jF$.

Let $O$ be the odd part of $t!$. Choose an arbitrarily large odd
$u=O(2m+1)$, and set $y=Fu$, $v=(u-1)/2$. The inequality
$v_2(F)\ge t>v_2(t!)$ gives $t!\mid y$. Choose $u$ large enough
that $y>F$ and exceeds every accumulator value. The beta moduli
$1+(i+1)y$ are pairwise coprime and odd. CRT gives a positive beta
code $x$, cofinal in its residue class.

Set $B=\binom{F+t}{t}$ and $N=B^x$. The two prime-power statements
assign compatible genuine rows at every prime factor of $B$.
The inherited precision-$N^2$ CRT construction supplies all nine
nonnegative representatives and range witnesses. Since $B$ is odd,
$N^2$ is odd, so their cofinal lifts also permit the parity and
dominance choices for the semantic anchor.

Now $H=3+6\sum a_j$ is positive, odd and divisible by three.
Thus $q=(BH)^x$ is positive, odd and divisible by three, and
$z=(q-3)/6$ is an integer. The standard canonical construction gives
the individual powers, Pell frame, normalized remainders and common
norm completion. The converse pooling theorem supplies its four new
coordinates. All source rows vanish. The exact projection and row
maps then give a zero of the final quartic.

Together with soundness and the inherited logical calibration, this
proves that the new arithmetic source, its pooled form, and its final
integer quartics have zeros exactly when there is a finite syntactic
ZFC refutation.


